\documentclass[11pt]{article}
\usepackage[a4paper,margin=27mm]{geometry}
\usepackage{amsmath,amssymb,amsthm,hyperref}
\newtheorem{lemma}{Lemma}
\newtheorem{proposition}[lemma]{Proposition}
\newcommand{\E}{\mathbb E}
\newcommand{\Pp}{\mathbb P}
\title{The own count with its conditional variance allocation retained}
\author{}
\date{28 October 2026}
\begin{document}
\maketitle

The difficult negative phases should not be bounded independently of the
incoming lower tail. Positive conditional noise is unfavorable for the
own count, whereas negative conditional noise is unfavorable for that
incoming tail. The elementary bound below retains this distinction. It
is a local estimate, not a completed count comparison above noise $1/50$.

Let $U_0,U_1$ be centered laws with
\[
 -A\le U_0\le a_0,\qquad -a_1\le U_1\le A,
 \qquad A>0,\quad 0\le a_0,a_1\le A.
\]
Write $v_j=\E U_j^2$, $\mu_j=\E U_j^3$. For an integer $n\ge1$ and
$0\le k\le n$, fix the legal counts
$N_0=n-k$ and $N_1=k$, and let $W$ be the sum of $N_j$ independent
variables of each type; all variables are independent. Set
$T_j=N_jv_j$ and $t=T_0+T_1$.
When marginal quantities are used, take $0<p<1$, $q=1-p$,
$B\sim\operatorname{Bernoulli}(p)$, and give $U$ the conditional law
$U_j$ given $B=j$. Then $s=\E U^2=qv_0+pv_1$ and $\lambda=ns$.

\begin{lemma}[A conditional fourth/third-moment bound]
\label{own:moment}
For every real $\kappa$,
\[
 \E W^4-\kappa\E W^3
 \le c_0(\kappa)T_0+c_1(\kappa)T_1+3t^2,
\]
where
\[
 c_0(\kappa)=\max\{A^2+\kappa A,a_0^2-\kappa a_0\},\qquad
 c_1(\kappa)=\max\{A^2-\kappa A,a_1^2+\kappa a_1\}.
\]
The right side can be reduced by $3(N_0v_0^2+N_1v_1^2)$ if that
correction is useful separately.
\end{lemma}
\begin{proof}
The convex function $x^2-\kappa x$ attains its maximum on a closed
interval at an endpoint. Multiplication by $x^2\ge0$ and integration
give $\E U_j^4-\kappa\E U_j^3\le c_j(\kappa)v_j$.
Centering and independence give
\[
 \E W^3=N_0\mu_0+N_1\mu_1,\qquad
 \E W^4=\sum\E U_i^4+3t^2-3\sum(\E U_i^2)^2.
\]
Substitution proves both statements.
\end{proof}

This improves the separate bounds $\E W^4\le A^2t+3t^2$ and
$|\E W^3|\le At$: each $c_j(\kappa)$ is at most
$A^2+A|\kappa|$. For nondegenerate support intervals the conditional
constants are asymptotically sharp.
For example, a centered two-point law on $\{-\varepsilon,A\}$ has
\[
 \frac{\E Z^4-\kappa\E Z^3}{\E Z^2}
 =A^2-A\varepsilon+\varepsilon^2-\kappa(A-\varepsilon)
 \longrightarrow A^2-\kappa A.
\]
The other three endpoints are obtained by reflection or by putting the
rare atom at an inner endpoint. This sharpness concerns the linear
per-summand moment bound, not an extremal convolution theorem.

\begin{proposition}[An own-count bound for any quartic minorant]
\label{own:quartic}
Suppose
\[
 h(x)=\frac{-x^4+b_3x^3+b_2x^2+b_1x+b_0}{M}
       \le\mathbf1_{(0,\infty)}(x),\qquad M>0.
\]
For $r\ge0$ define
\[
 \kappa=b_3-4r,\quad L=b_2+3b_3r-6r^2,\quad K=Mh(r).
\]
Then
\begin{equation}\label{own:lower}
 \Pp(W>-r)\ge
 \frac{Lt-3t^2-c_0(\kappa)T_0-c_1(\kappa)T_1+K}{M}.
\end{equation}
The bound includes atoms at the threshold: the event is strict.
\end{proposition}
\begin{proof}
Since $\E W=0$, expansion gives
$M\E h(W+r)=-\E W^4+\kappa\E W^3+Lt+K$.
Apply Lemma~\ref{own:moment} and the pointwise minorant.
\end{proof}

For the earlier quartic
\[
 h_\alpha(x)=\frac{-x^4+3\alpha^2x^2+2\alpha^3x}{5\alpha^4},
\]
the numerator factors as $-x(x+\alpha)^2(x-2\alpha)$; it is
nonpositive for $x\le0$ and its positive maximum is
$(9+6\sqrt3)\alpha^4/4<5\alpha^4$. Hence it is a valid minorant.
Here $\kappa=-4r$. If $a_0+a_1\le d\le A$ and
$0\le r\le(A-d)/4$, then
\[
 c_0=A^2-4Ar,\qquad c_1=A^2+4Ar.
\]
Indeed the first comparison is equivalent to $4r\le A-a_0$;
the second follows from $A^2+4Ar\ge a_1^2-4ra_1$.
Consequently
\begin{align}
 \Pp(W>-r)\ge\frac1{5\alpha^4}\{&
 (3\alpha^2-A^2-6r^2)t-3t^2
 +4Ar(T_0-T_1)\notag\\
 &-r^4+3\alpha^2r^2+2\alpha^3r\}.
 \label{own:symmetric}
\end{align}
Compared with the separate fourth/third-moment estimate, the gain is
exactly $8ArT_0/(5\alpha^4)$.

\paragraph{A material change at an admissible point.}
Take $A=3/5$, $\alpha=1/2$, $r=t=3/40$, and $T_1=0$.
After adding the normal-increment lower bound $-(21/20)r$, the new
lower bound is $55479/800000=.06934875$. The separate-moment bound
is $-13641/800000=-.01705125$. The gain is $.0864$.
These parameters are realized by an actual law: take $n=42$, $p=1/2$,
$k=21$, $U_1=0$, and the centered law $U_0\in\{-3/5,1/168\}$.
Then $ns=3/40$, $T_0=3/40$, $a_0+a_1=1/168<.072$ and the
corresponding raw threshold lies in $[-1.84,1.84]$.

\paragraph{Normal slope through phase $.1$.}
Put $v=pq$, $\ell=(v+s)^{-1/2}$, $\pi_k=\binom nkp^kq^{n-k}$,
and choose the phase by $np+H/\ell=k+u$. With $\delta=k-np$, the
normal increment is
\[
 f(u)=\frac{\Phi((\delta+u)/\sqrt{n(v+s)})-
                 \Phi(\delta/\sqrt{n(v+s)})}{\pi_k}.
\]
On $p\in[.4,.53]$, $n\ge42$, $ns\le.075$ and $|H|\le1.84$,
one has $\sqrt{pq+s}<.502$. For $|u|\le.1$ this gives
$|\delta|<1.84\cdot.502+.1=1.02368<21/20$.
The finite Stirling/logarithmic derivative estimate established in the
v10 comparison, with $D=21/20$ and $b=1/10$, gives
\[
 n|\log f'(z)|\le\tfrac83\lambda+
 \frac{2Db+b^2}{2(6/25)}+\frac{5D}{8}+\frac38+\frac{D^3}{42}.
\]
The exact right side at $\lambda=3/40$, divided by $42$, is less
than $1/21$. Thus $.95<f'(z)<1.05$ on the whole integration
segment. In particular $f(-r)\ge-(21/20)r$ for $0\le r\le.1$.
This invocation retains the finite Stirling proof and the fixed actual
count; it is not a fixed-law asymptotic assertion.

\section*{Combining the own count and incoming tails}
For fixed $n,p,k,\lambda>0$, put $q=1-p$ and parametrize the allocation
by $\eta=pv_1/s\in[0,1]$, where $s=\lambda/n$. Then
\[
 v_0=s(1-\eta)/q,\qquad v_1=s\eta/p,
 \quad T_0=N_0s(1-\eta)/q,\quad T_1=N_1s\eta/p.
\]
Fix the support caps and the quartic parameters independently of $\eta$.
The right side of \eqref{own:lower} is a concave function of $\eta$:
all terms are affine except $-3t^2$. Endpoint tests may use the larger
interval $[0,1]$ even if some endpoint allocations are infeasible for the
particular caps. They are then tests of a relaxation, not assumed laws.

For completeness, the needed incoming estimate can be proved without
an additional probability theorem. For $h>0$, $b>0$, define
\[
 B_h(b)=\frac{e^{hb}-1-hb}{b^2}
        =h^2\int_0^1(1-z)e^{hzb}\,dz.
\]
If $Z$ is centered and $Z\le b$, the same integral remainder gives
$\E e^{hZ}\le1+B_h(b)\E Z^2\le e^{B_h(b)\E Z^2}$.
Use $b=A$ for $-U_0$ and any fixed positive upper cap $b_1\ge a_1$
for $-U_1$. For each legal incoming count $i\ge1$,
\begin{equation}\label{own:chernoff}
 \Pp(W_{k+i}\le u-i)\le
 \exp\{h(u-i)+B_h(A)(N_0-i)v_0+B_h(b_1)(N_1+i)v_1\}.
\end{equation}
Every exponent is affine in $\eta$, so their sum is convex.
An infinite geometric extension is allowed if
\[
                  h> (s/p)B_h(b_1).
\]
Indeed the ratio is
$\exp\{-h-B_h(A)v_0+B_h(b_1)v_1\}<1$ uniformly on $[0,1]$.
The added terms are positive even when their count is no longer legal.
The infinite sum remains convex, being a pointwise limit of its convex
partial sums.

It follows that an own-count lower bound minus a fixed nonnegative
multiple of this incoming bound is concave in $\eta$. Its minimum is
at an allocation endpoint. This is an exact reduction to two scalar
estimates for each set of the other parameters. It does not itself
prove either endpoint estimate.

The tilt and quartic may depend on the other parameters, but must be
fixed while this allocation argument is applied. Taking a pointwise
minimum over several incoming bounds need not preserve convexity.
Likewise, retaining the favorable $3\sum v_i^2$ correction from
Lemma~\ref{own:moment} can destroy concavity in the allocation; dropping
it is deliberate here.

\section*{Retaining the cubic repair in the same inequality}
Suppose additionally $a_0<p$, $a_1<q$, and set
\[
 \ell=(pq+s)^{-1/2},\quad \rho=\E|\ell(B+U-p)|^3,
 \quad \rho_B=(1-2pq)/\sqrt{pq}.
\]
The signs are fixed, giving the exact identity
\[
 \rho=\ell^3\{pq(1-2pq)+3pq(v_0+v_1)+p\mu_1-q\mu_0\}.
\]
Let $m>0$, $C\ge0$, and assume $N_0,N_1>0$. Define
\[
 \gamma=C\ell^3M/m,\quad
 \kappa_0=\kappa-\gamma q/N_0,\quad
 \kappa_1=\kappa+\gamma p/N_1.
\]
The same proof, now combining the third moments before bounding them,
gives
\begin{align*}
 m\Pp(W>-r)+C(\rho-\rho_B)\ge{}&
 \frac mM\{Lt-3t^2-c_0(\kappa_0)T_0-c_1(\kappa_1)T_1+K\}\\
 &+C\bigl[\ell^3\{pq(1-2pq)+3pq(v_0+v_1)\}-\rho_B\bigr].
\end{align*}
For fixed $n,p,k,\lambda,m,C$ this is still concave in $\eta$.
In the comparison, $m=\sqrt n\,\pi_k$ is independent of $\eta$.
Outward replacements for $m$ must respect the sign of the complete
expression; one cannot replace it separately inside $\kappa_j$ without
checking that direction. This bound retains the same third moments in
both the own-count and marginal-moment terms.

\section*{An optional sharper bound at fixed summand variance}
The linear coefficient bound has a useful refinement, though it is not
needed for the endpoint reduction. If $Z\in[-b,A]$ is centered with
variance $v$ and
\[
                   A-b-2v/A-\kappa\ge0,
\]
then
\[
 \E(Z^4-\kappa Z^3)
 \le v\{A^2-v+v^2/A^2-\kappa(A-v/A)\}.
\]
To prove this, subtract
$(x-A)(x+v/A)^2(x+A-2v/A-\kappa)$ from $x^4-\kappa x^3$.
The resulting polynomial is quadratic and majorizes the original on
$[-b,A]$. Its expectation depends only on the mean and variance; evaluate
it on the centered law supported at $-v/A,A$. This law has the required
variance, and $v/A\le b$ follows from $v\le Ab$. At both its atoms
the subtraction vanishes. The case $v=0$ is immediate. Reflection
gives the analogous negative-endpoint estimate. This refinement should
not be inserted into the allocation argument without rechecking its
curvature.

\paragraph{Status.}
The lemmas are complete local deductions. Exact rational checks cover
the normal-slope constants, moment identities, support endpoints,
actual conditional laws and illustrative improvement. No complete
comparison on $.02<\lambda\le.075$ is claimed. The remaining work is
to prove the two allocation-endpoint comparisons with one common
phase-dependent choice of tilt and own-count bound.
\end{document}
